\chapter{Удиртгал}

\section{Сэдвийн үндэслэл}
Орчин үед байгууллага, компаниуд хэрэглэгчдийнхээ санал сэтгэгдлийг сонсох, тренд болж буй сэдвүүдийг илрүүлэх зорилгоор сошиал мониторингийн системийг өргөнөөр ашиглах болсон. Одоогийн байдлаар Sumbee.mn зэрэг системүүд нь ихэвчлэн түлхүүр үг дээр тулгуурласан тоон болон эерэг/сөрөг (sentiment) анализ хийдэг бөгөөд агуулгын түвшинд задлан шинжлэх, илүү бодитой, бүтэцчилсэн дүгнэлт гаргах боломж дутмаг байна. Мөн маш их хэмжээний өгөгдөл дээр LLM шууд ашиглах нь тооцооллын зардал маш их шаардах ба удаан (latency өндөр) байдаг. Иймд Монгол хэл, тэр дундаа сошиал өгөгдөл дээр үр дүнтэй ажиллаж чадах NLP (Байгалийн хэлний боловсруулалт)-д суурилсан шийдэл боловсруулах шаардлага тулгарч байна.

\section{Зорилго болон Зорилт}
Энэхүү судалгааны ажлын зорилго нь сошиал сүлжээний бүтэцлэгдээгүй текст өгөгдлөөс далд сэдвүүд (Topic) болон нэрлэсэн объектуудыг (NER) автоматаар илрүүлж, тэдгээрийн хамаарлын сүлжээг (Network Analysis) үүсгэх инференс пайплайн зохиож хэрэгжүүлэхэд оршино. Үүний тулд дараах зорилтуудыг дэвшүүллээ:
\begin{itemize}
    \item Одоо байгаа NLP болон Topic modeling (LDA, NMF, BERTopic) аргуудыг судлах, харьцуулах.
    \item Монгол хэлний pretrained BERT загварыг ашиглан текст өгөгдлийг векторжуулах (Embedding).
    \item HuggingFace-ийн pretrained NER загварыг (жишээлбэл Davlan/bert-base-multilingual-cased-ner-hrl) турших.
    \item Сошиал өгөгдөл дээр сэдэв илрүүлэх болон нэрлэсэн объект таних хосолсон архитектур үүсгэх.
    \item Илрүүлсэн объектуудын (хүн, байгууллага, газар) хоорондын хамаарлыг Network Analysis ашиглан график дүрслэлд оруулах.
\end{itemize}

\section{Ажлын бүтэц}
Энэхүү тайлан нь нийт 5 бүлгээс бүрдэнэ. Нэгдүгээр бүлэгт удиртгал, хоёрдугаар бүлэгт текст боловсруулалт болон сэдэв загварчлалын онолын суурь судалгаа, гуравдугаар бүлэгт системийн шинжилгээ болон архитектур зохиомж, дөрөвдүгээр бүлэгт хэрэгжүүлэлт болон туршилтын үр дүн, тавдугаар бүлэгт дүгнэлтийг тус тус тусгасан.
